\documentclass[11pt]{article}

% "review"  -> anónimo, con números de línea (para envío a revisión)
% "final"   -> versión final (nombres de autor, sin números de página)
% "preprint"-> no anónimo, con números de página  <-- usamos esta
\usepackage[preprint]{acl}

% Paquetes estándar
\usepackage{times}
\usepackage{latexsym}
\usepackage[T1]{fontenc}
\usepackage[utf8]{inputenc}
\usepackage{microtype}
\usepackage{inconsolata}
\usepackage{graphicx}

% Español. es-noshorthands evita que babel active atajos (" < > ~ .) que
% chocarían con Tabla~\ref, math, etc.
\usepackage[spanish,es-noshorthands]{babel}
\renewcommand{\abstractname}{Resumen}
\addto\captionsspanish{\renewcommand{\tablename}{Tabla}}

% Para las tablas del trabajo
\usepackage{booktabs}
\usepackage{amsmath}
\usepackage{float}

\title{MiroFish: Simulación Social Generativa y Evaluación de \\ Sistemas Multi-Agente en Escenarios Complejos}

\author{
  \textbf{Andrés Bravo}, \textbf{Joaquín Cadeiras}, \textbf{Bruno De Cruz}, \\
  \textbf{Lucas Ercolano}, \textbf{Eliana Ostrovsky} \\
  Universidad de San Andrés \\
  \small{Procesamiento del Lenguaje Natural}
}

\begin{document}
\maketitle

\begin{abstract}
Los sistemas de simulación social generativa, como MiroFish, pronostican eventos haciendo deliberar a poblaciones de agentes basados en modelos de lenguaje, pero típicamente ejecutan a todos los agentes sobre un mismo modelo. Proponemos poblar la simulación con modelos heterogéneos (Llama-3.3-70B, Gemma-3-27B y Qwen3-8B) bajo la hipótesis de que ello reduciría el sesgo y el costo. Como no existía un punto de comparación, construimos un \textit{benchmark} de tres casos —económico, político y deportivo— con control de fugas de información, y evaluamos los modelos aislados frente a la red heterogénea bajo los mismos regímenes de estrés. Encontramos que, para el pronóstico puro, un modelo aislado y calibrado es más preciso y más barato, mientras que la red heterogénea, sin estímulos externos, sufre un colapso semántico que hunde su precisión; el ahorro de costo real provino de una deduplicación semántica del grafo de conocimiento. La heterogeneidad no auto-regula el sesgo, pero sí lo vuelve permeable a inyecciones exógenas. Concluimos que el valor de la red multi-modelo no está en el \textit{forecasting} sino como laboratorio para estudiar dinámicas socio-técnicas de persuasión y desinformación.
\end{abstract}

\section{Introducción}
\label{sec:intro}

Los grandes modelos de lenguaje (LLM) hicieron viable simular poblaciones de agentes que conversan, se influyen y toman decisiones, abriendo la puerta a emplear estas simulaciones como instrumento de pronóstico de eventos reales \citep{park2023generative, yang2024oasis}. MiroFish \citep{mirofish2025} es un sistema de esta familia: a partir de un corpus semilla instancia una población de agentes en un foro social, los hace deliberar y destila un pronóstico a partir del consenso emergente. La premisa es que, en eventos cuyo desenlace depende de la reacción colectiva, la dinámica social simulada aporta señal más allá de la extrapolación puramente estadística.

Ahora bien, estos entornos suelen ejecutar a todos los agentes sobre un mismo LLM. Una población algorítmicamente homogénea tiende a reproducir las patologías de su único modelo subyacente —cámaras de eco, rigidez de pre-entrenamiento, baja diversidad léxica— y, cuando ese modelo es grande, el costo de cómputo de una simulación completa se vuelve considerable. Postulamos que poblar la simulación con modelos heterogéneos —de distinto tamaño, corpus de pre-entrenamiento y alineación— podría mitigar ambos problemas: la fricción entre arquitecturas divergentes reduciría el sesgo, y la incorporación de modelos más pequeños abarataría la inferencia.

Contrastar esta hipótesis exigía un punto de comparación inexistente: MiroFish no contaba con un \textit{benchmark} que permitiera medir si una modificación mejora o degrada el sistema. Por eso construimos uno propio, con tres casos de distinta temática (económico, político y deportivo) bajo un protocolo estricto de control de fugas de información, y evaluamos primero los modelos en aislamiento (\textit{baseline}) y luego la red multi-modelo, replicando en ambos los mismos regímenes de estrés.

Este trabajo aporta: (i) un mecanismo de ruteo heterogéneo de modelos por agente, con telemetría de costo por llamada; (ii) una deduplicación semántica del grafo de conocimiento que reduce el cómputo sin perder información; (iii) un \textit{benchmark} reproducible de tres casos con \textit{ground truth} aislado; y (iv) un análisis empírico, a lo largo de varias líneas experimentales, de cómo distintas variables afectan la simulación y de en qué difiere el comportamiento de una población homogénea frente a una heterogénea.

\section{Trabajo Relacionado}
\label{sec:related}

\paragraph{Simulación social generativa.} La idea de instanciar agentes basados en LLM que interactúan de forma creíble fue popularizada por los \textit{generative agents} de \citet{park2023generative}, y llevada a escala poblacional por OASIS \citep{yang2024oasis}, que simula hasta un millón de usuarios sobre clones de Twitter y Reddit. MiroFish \citep{mirofish2025} reutiliza este tipo de entorno pero lo orienta al pronóstico de eventos. A diferencia de estos trabajos, que instancian una población homogénea sobre un único modelo, nuestro foco es la heterogeneidad de modelos dentro de una misma simulación y su efecto sobre el sesgo, el costo y la calidad predictiva.

\paragraph{Degeneración y diversidad.} La degradación de sistemas generativos realimentados con su propia salida (\textit{model collapse}) fue caracterizada por \citet{shumailov2024collapse}; observamos un fenómeno análogo cuando la red delibera en circuito cerrado sin entropía exógena. Para cuantificar la variedad lingüística de cada modelo —de forma independiente de su puntería predictiva— empleamos métricas estándar de diversidad: distinct-$n$ \citep{li2016diversity}, Self-BLEU \citep{zhu2018texygen} y Vendi Score \citep{friedman2023vendi}. Por último, dado que las asimetrías de comportamiento entre modelos (proactividad, permeabilidad, cortesía evasiva) se originan en parte en sus esquemas de alineación \citep{ouyang2022instructgpt}, también interpretamos las dinámicas inter-modelo a esa luz.

\section{Metodología}
\label{sec:metodologia}

\subsection{Arquitectura del sistema}

MiroFish \citep{mirofish2025} es un motor de predicción que estima el desenlace de un evento simulando la conversación de una población de agentes en un foro social. A partir de un corpus semilla (noticias, informes, encuestas) el sistema construye un grafo de conocimiento, instancia agentes con perfiles y memoria propios, y los hace interactuar durante un número configurable de rondas; el consenso emergente se destila luego en una predicción. La capa de interacción social reutiliza el entorno OASIS \citep{yang2024oasis}, que clona plataformas tipo Twitter y Reddit e incorpora sistemas de recomendación y un repertorio de acciones (publicar, comentar, seguir, repostear, entre otras).

Sobre esta base introducimos las dos modificaciones que constituyen el aporte técnico de este trabajo:

\paragraph{Ruteo heterogéneo de modelos.} Por defecto, todos los agentes de una simulación corren sobre un mismo LLM. Implementamos un enrutador (\texttt{model\_router.py}) que asigna modelos distintos a distintos agentes mediante un mapa declarativo en YAML, junto con una capa de telemetría por llamada (\texttt{llm\_telemetry.py}) que registra tokens, latencia y costo de cada acción. Esto permite poblar un mismo foro con arquitecturas de pre-entrenamiento y alineación divergentes, y atribuir el costo de cómputo a cada modelo.

\paragraph{Deduplicación semántica del grafo.} Durante la ingesta observamos la generación de nodos redundantes o casi idénticos que aumentan el costo y el tiempo de cómputo sin aportar información nueva. Añadimos un paso que fusiona los nodos semánticamente equivalentes antes de insertarlos en el grafo (Neo4j, con \textit{embeddings} \texttt{bge-m3}), reduciendo su tamaño. Su impacto se cuantifica en la Sección~\ref{sec:evaluacion}.

Todas las simulaciones se ejecutan en entornos herméticos: los agentes no acceden a Internet y el corpus semilla respeta un corte temporal estricto anterior al evento, evitando fugas de información (\textit{data leakage}). El resultado real se mantiene aislado del input y solo se emplea en la evaluación.

\subsection{Casos de estudio}
\label{subsec:casos}

Al no existir un \textit{benchmark} previo para el sistema, construimos uno propio con tres casos de distinta temática (Tabla~\ref{tab:casos}): económico, político y deportivo. Cada caso plantea un evento cuyo desenlace es posterior al corte documental y verificable de forma objetiva, de modo que el \textit{benchmark} cubre dominios con dinámicas de información heterogéneas manteniendo una evaluación no ambigua. Al representar cada temática con un único evento, el objetivo es la amplitud de dominios y no una medición estadísticamente concluyente dentro de cada uno; una caracterización más fina de las fortalezas y debilidades del sistema por dominio queda como trabajo futuro.

\paragraph{Económico: inflación del IPC (Argentina, 2025).} La variación mensual del Índice de Precios al Consumidor que publica el INDEC. Fijamos el corte documental el 31 de enero de 2025 y pedimos pronosticar el IPC mensual en cuatro horizontes crecientes ($\Delta$1 a $\Delta$4: febrero, abril, julio y diciembre de 2025), para evaluar si el sistema captura la trayectoria de desinflación del programa económico entonces vigente. Es el caso más cuantitativo y cerrado: el desenlace es un valor oficial no ambiguo. El \textit{ground truth} muestra una desinflación real pero no lineal —con un rebote a mitad de año— y una inflación acumulada anual en torno al 31.5\%.

\paragraph{Político: balotaje presidencial (Bolivia, 2025).} La segunda vuelta del 19 de octubre de 2025 entre Rodrigo Paz y Jorge ``Tuto'' Quiroga, en un ciclo marcado por la crisis económica y el colapso electoral del MAS tras casi dos décadas en el poder. A partir de encuestas, resultados de primera vuelta y contexto socioeconómico previos al balotaje, el sistema debe anticipar el ganador y su porcentaje de voto. Es el caso más abierto y polarizado. El resultado real fue el triunfo de Paz con 54.53\% frente al 45.47\% de Quiroga (margen de 9.06 puntos).

\paragraph{Deportivo: final de la Copa América (2024).} La final del 14 de julio de 2024 entre Argentina y Colombia, disputada en Miami. Con un corte documental el día previo (13 de julio), el sistema debe predecir el ganador a partir de previas deportivas, estado de forma e historial. Es un caso simple y auto-verificable, de horizonte muy corto ($\sim$1 día), que opera como prueba de sanidad del protocolo temporal. El desenlace real fue la victoria de Argentina por 1--0 en tiempo suplementario.

\begin{table}[t]
\centering
\small
\begin{tabular}{@{}llll@{}}
\toprule
\textbf{Tema} & \textbf{Caso} & \textbf{Predicción} & \textbf{Métrica} \\
\midrule
Econ. & IPC Arg. 2025 & Inflación ($\Delta$1--$\Delta$4) & MAE (\%) \\
Polít. & Bolivia 2025 & Ganador / voto & MAE (voto) \\
Dep. & Copa Am. 2024 & Ganador & Acierto \\
\bottomrule
\end{tabular}
\caption{Los tres casos del \textit{benchmark}: inflación mensual (INDEC), balotaje presidencial (Paz vs.\ Quiroga) y la final Argentina--Colombia. En todos, el evento es posterior al corte documental y el \textit{ground truth} se mantiene aislado del input.}
\label{tab:casos}
\end{table}

\subsection{Modelos evaluados}

Evaluamos tres LLM de código abierto seleccionados para desacoplar dos variables habitualmente correlacionadas: el tamaño del modelo y la novedad de su fecha de corte. Elegimos deliberadamente una relación invertida entre ambas: \texttt{Llama-3.3-70B} \citep{llama33} es el mayor pero el de corte de conocimiento más antiguo (diciembre de 2023), mientras que \texttt{Qwen3-8B} \citep{qwen3} es el menor y el de corte más reciente (marzo de 2025); \texttt{Gemma-3-27B} \citep{gemma3}, con corte en agosto de 2024, ocupa una posición intermedia en ambos ejes.\footnote{Qwen no publica una fecha de corte oficial para Qwen3; la estimamos en marzo de 2025 sondeando el modelo con una batería de preguntas sobre eventos datados. Debe tomarse como aproximación, no como dato oficial.} Este diseño permite indagar si las patologías que se observen son atribuibles a la escala del modelo o a la cobertura temporal de su pre-entrenamiento.

\subsection{Métricas}

Distinguimos dos familias. Para el \textbf{poder predictivo} (\textit{forecasting}) usamos el Error Absoluto Medio (MAE) respecto del \textit{ground truth}, en la unidad natural de cada caso: puntos porcentuales de inflación (IPC), puntos de voto (Bolivia) y acierto binario del ganador (Copa América). Para el \textbf{comportamiento intrínseco} de los modelos —independiente de su puntería— empleamos métricas agnósticas de diversidad y estabilidad: Self-BLEU y distinct-2 \citep{li2016diversity,zhu2018texygen}, Vendi Score \citep{friedman2023vendi}, entropía categórica y deriva temporal intra-persona, esta última medida con entrevistas paramétricas a los agentes en distintos puntos de la simulación (\textit{Checkpoints Interview}). La capa de telemetría complementa ambas familias registrando el consumo de tokens, el costo en USD y la adherencia al formato de salida (JSON).

% ============================================================================
% RESULTADOS: primero TODO el baseline, después TODO el multimodelo (pedido del
% equipo). Cada "línea" (variable estudiada) es una subsubsección separada.
% ============================================================================
\section{Evaluación}
\label{sec:evaluacion}

\subsection{Baseline: modelos aislados (Single-Agent)}
\label{subsec:baseline}

Como línea de base ejecutamos cada modelo en aislamiento —toda la población sobre un mismo LLM— y sometimos el sistema a una batería de regímenes de estrés. Salvo indicación contraria, usamos \texttt{Llama-3.3-70B} como modelo ancla por ser el de mayor precisión en el caso cuantitativo.

\subsubsection{Escalamiento topológico (rondas y densidad)}

Variamos la profundidad del debate (número de rondas $R$) y la densidad de conexión del grafo de agentes ($D$). El efecto resultó opuesto según la naturaleza del caso (Tabla~\ref{tab:topologia}). En un entorno cuantitativo y cerrado (IPC), más deliberación mejora la inferencia: escalar a $R=80$, $D=2$ redujo el MAE a un óptimo de 1.84\%. En cambio, en un entorno abierto y polarizado (balotaje de Bolivia), el mismo escalamiento retroalimenta las asunciones del modelo y forma una cámara de eco (\textit{herd behavior}): el error crece de 7.60 a 10.0 puntos de voto al pasar de 10 a 40 rondas.

\begin{table}[t]
\centering
\small
\begin{tabular}{@{}llc@{}}
\toprule
\textbf{Entorno} & \textbf{$R$/$D$} & \textbf{MAE} \\
\midrule
IPC (cerrado) & 10/2 & 3.12\% \\
IPC (cerrado) & 40/1 & 2.45\% \\
IPC (cerrado) & 40/2 & 2.31\% \\
\textbf{IPC (cerrado)} & \textbf{80/2} & \textbf{1.84\%} \\
\midrule
Bolivia (abierto) & 10/2 & 7.60 \\
\textbf{Bolivia (abierto)} & \textbf{40/2} & \textbf{10.0} \\
\bottomrule
\end{tabular}
\caption{Escalamiento topológico con \texttt{Llama-3.3-70B}. Más rondas mejoran el caso cuantitativo cerrado (IPC, MAE en \%) pero desencadenan \textit{herd behavior} en el abierto (Bolivia, MAE en puntos de voto).}
\label{tab:topologia}
\end{table}

\subsubsection{Sesgos por arquitectura}

Comparamos el rendimiento de las tres arquitecturas en el caso cuantitativo (Tabla~\ref{tab:arquitecturas}). \texttt{Llama-3.3-70B} alcanza el mejor MAE medio (1.85\%, $\sigma=0.05$ sobre réplicas de su configuración óptima $R=80$, $D=2$), seguido de \texttt{Gemma-3-27B} (2.67\%) y \texttt{Qwen3-8B} (3.45\%) en la configuración de control ($R=40$, $D=2$). En escenarios sociales abiertos, cada arquitectura manifiesta una patología distinta: Llama forma una burbuja epistemológica en el caso electoral; Gemma exhibe una impermeabilidad factual absoluta en el deportivo —ignora las estadísticas contrarias a su sesgo, anuladas por el peso de su pre-entrenamiento—; y Qwen cae en un consenso neutro y pasivo, aunque es el más permeable a inyecciones de datos sorpresivos. Estas asimetrías motivan directamente el diseño multi-modelo (§\ref{subsec:multimodelo}).

\begin{table}[t]
\centering
\small
\begin{tabular}{@{}lcl@{}}
\toprule
\textbf{Modelo} & \textbf{MAE (IPC)} & \textbf{Patología (abierto)} \\
\midrule
Llama-3.3-70B & 1.85\%$^{*}$ & Burbuja epistemológica \\
Gemma-3-27B & 2.67\% & Terquedad (ignora datos) \\
Qwen3-8B & 3.45\% & Consenso neutro \\
\bottomrule
\end{tabular}
\caption{Cruce de arquitecturas: precisión en el caso cuantitativo (IPC) y patología dominante en escenarios abiertos. $^{*}$Llama medido en réplicas de control (R80-D2).}
\label{tab:arquitecturas}
\end{table}

\subsubsection{Deduplicación de nodos: costo y precisión}

Durante la construcción del grafo detectamos una proliferación de nodos redundantes o casi idénticos, que aumentan el costo y el tiempo de cómputo. Añadimos un paso que verifica la similitud semántica antes de instanciar un nodo y descarta los nodos sin conectividad. En el caso IPC, esta optimización mejoró la precisión respecto de la ingesta manual sin deduplicación (MAE de 2.31\% a 1.475\%, la mejor marca del proyecto) además de reducir el número de nodos y, con ello, el costo. Al tratarse de una única prueba no podemos afirmar que la deduplicación mejore el pronóstico en general; pero en este caso redujo simultáneamente el cómputo y el error.
% Nota: MAE 2.31%->1.475% CONFIRMADO en s3_final_report.md (origin/backtesting-baseline).
% El "-10% de nodos" del equipo NO aparece literal en el artefacto; por eso acá va
% cualitativo ("reducir el número de nodos"). Si se consigue el conteo exacto, agregarlo.

\subsubsection{Permeabilidad e ingesta autónoma (Deep Search)}

Evaluamos dotar al sistema de recolección de contexto totalmente autónoma (\textit{Deep Search}) con cortes temporales estrictos. Lejos de mejorar, la precisión se degradó: en ausencia de anclas documentales estructuradas, la atención de los agentes se dispersa y las proyecciones pierden consistencia. La ingesta curada del corpus semilla resultó preferible a la recolección autónoma.

\subsubsection{Sensibilidad a la ventana temporal (\texorpdfstring{$\Delta$1--$\Delta$4}{Delta 1-4})}

Desplazando la ventana temporal de los datos semilla medimos la degradación del pronóstico en el caso IPC, con un patrón no lineal. A corto plazo ($\Delta 1$) el sistema es muy preciso (error de unos 0.1 puntos porcentuales), anclado en valores nominales presentes en el corpus. El mayor fallo ocurre a medio-corto plazo ($\Delta 2$): sin señales sobre \textit{shocks} transitorios futuros, el modelo extrapola linealmente la tendencia inicial y subestima los rebotes macroeconómicos. A medio plazo ($\Delta 3$) la precisión se recupera para la tendencia de fondo. El motor es, entonces, robusto para la dirección dominante pero ciego a las fluctuaciones transitorias si la ventana no se actualiza.

\subsubsection{Robustez ante ruido (inyecciones exógenas)}

Inyectamos ruido deliberado en el entorno (p.\ ej.\ especulación cambiaria) y variamos su momento de entrada. Bajo alta profundidad ($R=80$, $D=2$), el modelo asimiló la especulación social contenida en el ruido y, paradójicamente, ancló un escenario recesivo con gran precisión (MAE de 0.9875\% en esa variante). El momento resultó crítico: datos correctos inyectados a mitad del debate suelen ser ignorados por arquitecturas con fuerte sesgo de pre-entrenamiento (Gemma), lo que confirma que la asimilación depende del balance entre el ruido externo y el arraigo de las creencias ya formadas.

\subsubsection{Diversidad léxica y estabilidad (evaluación intrínseca)}

Con independencia de la puntería predictiva, medimos el funcionamiento intrínseco de cada modelo con métricas algorítmicas agnósticas (entropía categórica, Self-BLEU, Vendi Score y deriva euclidiana). Los resultados exponen asimetrías severas en el comportamiento generativo (Tabla~\ref{tab:intrinseca}).

\begin{enumerate}
    \item \textbf{Composición y proactividad}: Llama es marcadamente reactivo —por cada 10 publicaciones originales genera 57 comentarios a terceros—, mientras que Gemma muestra un perfil proactivo e independiente (18 publicaciones, 16 comentarios).
    \item \textbf{Entropía léxica}: sobre publicaciones originales, Gemma posee una diversidad semántica marcadamente mayor (Self-BLEU 0.264, distinct-2 0.752), mientras que Llama recicla estructuras léxicas de forma sistemática (Self-BLEU 0.412).
    \item \textbf{Deriva temporal (\textit{intra-persona})}: mediante entrevistas paramétricas al inicio, la mitad y el cierre de la simulación, Llama exhibió mayor inestabilidad (\textit{endpoint-distance} 0.266) que Qwen (0.225); las personas simuladas por Llama tienden a mutar sus posturas si la simulación se prolonga.
\end{enumerate}

\begin{table*}[t]
\centering
\small
\begin{tabular}{@{}lccccc@{}}
\toprule
\textbf{Modelo} & \textbf{Posts} & \textbf{Coment.} & \textbf{Self-BLEU}$\downarrow$ & \textbf{distinct-2}$\uparrow$ & \textbf{Deriva} \\
\midrule
Llama-3.3-70B & 10 & 57 & 0.412 & 0.614 & 0.266 \\
\textbf{Gemma-3-27B} & 18 & 16 & \textbf{0.264} & \textbf{0.752} & N/A \\
Qwen3-8B & --- & --- & --- & --- & \textbf{0.225} \\
\bottomrule
\end{tabular}
\caption{Evaluación intrínseca. Gemma lidera en diversidad léxica y proactividad relacional; Qwen, en estabilidad ideológica (menor deriva temporal). Las celdas sin valor (---, N/A) corresponden a métricas no computadas para ese modelo en esta fase.}
\label{tab:intrinseca}
\end{table*}

\subsection{Multi-modelo: red heterogénea}
\label{subsec:multimodelo}

Motivados por las asimetrías del baseline, poblamos una misma simulación con los tres modelos y replicamos los mismos regímenes de estrés. Nuestra hipótesis era que la fricción entre arquitecturas divergentes funcionaría como auto-regulación y reduciría el sesgo.

\subsubsection{Escalamiento de rondas y colapso semántico}

En el baseline, escalar las rondas con un solo modelo (Llama) inflaba la cámara de eco hasta un MAE de 10.0 en el balotaje de Bolivia. Al replicar la prueba con la red heterogénea y sin estímulos externos, la hipótesis de auto-regulación se refutó de forma contundente: en lugar de corregir hacia la verdad, el sistema sufrió un \textbf{colapso semántico} (\textit{model collapse}). Forzados a dialogar sin nueva entropía, los agentes cayeron en un \textit{deadlock} dialéctico —la hiper-reactividad de Llama frente a la proactividad de Gemma, moduladas por sus esquemas de alineación— y el foro se llenó de cortesía vacía (p.\ ej.\ \textit{``I completely agree with the importance of considering the candidates' policies...''}), sin que ningún modelo tomara postura.

Como resultado, el sistema falló en extraer un consenso válido: el MAE se disparó a 66.667 debido a un abstencionismo del 100\% (clasificado como indecisos/otros), muy por encima del error del modelo aislado (Tabla~\ref{tab:collapse}). En cautiverio y sin nueva entropía, la diversidad de modelos degeneró en pasividad absoluta.

\begin{table}[t]
\centering
\small
\begin{tabular}{@{}llc@{}}
\toprule
\textbf{Arquitectura} & \textbf{Predicción} & \textbf{MAE} \\
\midrule
Baseline (Llama), R10 & Quiroga & 7.687 \\
Baseline (Llama), R40 & Quiroga & 10.000 \\
Multi-modelo, R10 & Indecisos & 66.667$^{*}$ \\
Multi-modelo, R40 & Indecisos & 66.667$^{*}$ \\
\bottomrule
\end{tabular}
\caption{Comparativa de MAE en el balotaje de Bolivia sin inyecciones externas. $^{*}$Denota colapso semántico total (abstencionismo del 100\%).}
\label{tab:collapse}
\end{table}

\subsubsection{Inyección dinámica (variación de momento)}

En el baseline, Gemma en solitario se mostró impenetrable ante las estadísticas a favor de Colombia en la Copa América y mantuvo obstinadamente su voto por Argentina. Al transicionar a la red multi-modelo, inyectamos la misma estadística a mitad del debate, con Llama-3.3 como vector de introducción. Sometida a la presión argumentativa de sus pares, Gemma fracturó su rigidez de pre-entrenamiento: el rastreo de la base de datos documentó el viraje cognitivo, con Gemma respondiendo \textit{``An excellent analysis! Colombia has a very real chance to shine...''}. Esto demuestra que la influencia inter-modelo es factible: la presión de los pares logró quebrar una rigidez que ningún estímulo conseguía mover en aislamiento. Ahora bien, ese mismo viraje empujó la predicción hacia Colombia, que finalmente perdió la final ante Argentina: la maleabilidad inducida por los pares no equivale a acierto.

\subsubsection{Shock de ruido tardío (entropía exógena)}

Sabiendo que la red colapsaba en neutralidad, inyectamos una encuesta falsa tardía a mitad del debate inactivo. El shock de entropía desestabilizó el estado de cortesía: Qwen, con la permeabilidad que ya mostraba en el baseline, absorbió el estímulo de inmediato, lo tradujo en retórica agresiva e impulsó la polarización del foro, rescatando al sistema del colapso. Los LLM heterogéneos requieren señales externas continuas para sostener una deliberación sana.

\subsubsection{Superación de la degradación temporal}

En el baseline, alterar de forma estática el período de los datos semilla producía la curva de degradación temporal (fallo en $\Delta 2$ por extrapolación lineal). En la red multi-modelo abandonamos el corte estático y alimentamos a los agentes con flujos iterativos que simulan el paso del tiempo. El ecosistema heterogéneo —sometido a fricción dialéctica constante— procesó las actualizaciones sin caer en la extrapolación lineal ciega del modelo aislado: la diversidad arquitectónica actuó como amortiguador predictivo, estabilizando la inferencia a corto y medio plazo mientras el flujo de inyección no se interrumpiera.

\section{Conclusiones}
\label{sec:conclusiones}
El contraste entre el \textit{baseline} y la arquitectura multi-modelo nos permite cristalizar tres conclusiones sobre la simulación social generativa.

\paragraph{Costo-beneficio en \textit{forecasting}.} Si el objetivo exclusivo es la predicción, los sistemas multi-modelo resultan computacionalmente prohibitivos e ineficientes: el consumo supera los 230.000 \textit{tokens} por debate y la propensión al colapso semántico deja su precisión por debajo de la de un modelo individual calibrado. La reducción de costo efectiva provino, en cambio, de la deduplicación semántica del grafo.

\paragraph{Fallo estructural del debate aislado.} La diversidad de modelos no es un antídoto automático contra las burbujas. En ausencia de un flujo de información exógeno, los modelos interactuando entre sí convergen hacia una neutralidad evasiva, producto de \textit{deadlocks} entre asimetrías conductuales severas y sus alineamientos éticos (RLHF).

\paragraph{El verdadero valor como laboratorio socio-técnico.} MiroFish alcanza su máximo cuando la red multi-modelo se somete a inyecciones exógenas paramétricas. Bajo esas condiciones, la arquitectura heterogénea es el mecanismo que validamos para simular persuasión real, absorción de desinformación y fractura de rigideces ideológicas, volviéndose una herramienta valiosa para estudiar la propagación de narrativas a escala poblacional.

% La sección de Limitaciones es OBLIGATORIA en ACL y NO cuenta para el límite
% de 8 páginas.
\section*{Limitaciones}

Este trabajo tiene un alcance exploratorio y sus conclusiones deben leerse con varias salvedades.

\paragraph{Cobertura muestral.} El \textit{benchmark} usa un único evento por temática (económico, político y deportivo). Esto permite abarcar dominios diversos, pero impide afirmaciones estadísticamente concluyentes sobre las fortalezas o debilidades del sistema en cada uno; distinguir un efecto de dominio del ruido de un caso particular requeriría varios eventos por temática.

\paragraph{Control de varianza parcial.} Solo algunas mediciones se repitieron con réplicas y control de varianza (la precisión del IPC y el análisis intrínseco de diversidad y deriva). Otros resultados —el escalamiento topológico, la inyección de ruido y el colapso multi-modelo— provienen de corridas únicas, por lo que sus valores puntuales deben tomarse como indicativos y no como estimaciones robustas.

\paragraph{Cobertura de modelos.} La comparación intrínseca se realizó entre Gemma-3 y Llama-3.3; Qwen3 no pudo incluirse por falta de acceso al modelo durante esa fase. Además, evaluamos una única terna de modelos en un rango de tamaño acotado (8B--70B) y una sola configuración de mezcla. No podemos garantizar que las conclusiones sean agnósticas al tamaño ni a la elección concreta de modelos; explorar otras combinaciones y escalas mayores es trabajo pendiente.

\paragraph{Métricas de evaluación.} El MAE de participación electoral se calcula a partir de la extracción, mediante expresiones regulares, de porcentajes en el informe generado; una predicción mal formada o evasiva puede no ser detectada. El MAE se expresa además en unidades distintas según el caso (puntos porcentuales de inflación frente a puntos de voto), por lo que los valores no son comparables entre casos.

\paragraph{Reproducibilidad.} Las simulaciones dependen de modelos servidos por terceros y de estimaciones de costo obtenidas por telemetría, no de un perfilado exhaustivo. La deduplicación de nodos, por su parte, se validó en un único caso, sin cuantificar de forma independiente la reducción exacta del tamaño del grafo.

\bibliography{custom}

\end{document}
